% SPDX-License-Identifier: GPL-3.0-or-later
% Copyright (C) 2026 Alessandro Greco
%% EPUB 无障碍替代文本。
\altfig{三种架构对比。（a）Tor Browser 只把浏览器自身的流量交给 tor，宿主机的其他应用程序不会自动受到保护，仍可能直接连接外网。虚线边框只圈出浏览器和 tor。（b）Tails 是从 USB 启动、主要运行在内存中的实时操作系统，通常以防火墙阻止应用程序绕过 Tor 直接联网；但也有刻意保留的例外，例如用于登录 captive portal、不会经过 Tor 的 Unsafe Browser。（c）Whonix 由运行应用的 Workstation 虚拟机和运行 tor 的 Gateway 虚拟机组成，二者通过隔离的内部虚拟网络连接。若宿主机或虚拟机管理程序被攻破，这种隔离仍可能失效。}
% 示意：浏览器、使用防火墙的 Live 系统、两台虚拟机
\begin{tikzpicture}[font=\sffamily\scriptsize,
  box/.style={draw,rounded corners=2pt,align=center,minimum width=2.2cm,minimum height=0.7cm},
  conf/.style={draw=torviola,dashed,rounded corners=4pt,inner sep=5pt}]
  % Tor Browser
  \node[box] (tb) at (0,1.6) {Tor Browser};
  \node[box,fill=torviola!15] (t1) at (0,0.4) {tor};
  \node[box,fill=torgrigio!15] (a1) at (0,-0.8) {其他应用};
  \draw[->] (tb) -- node[right] {SOCKS} (t1);
  \draw[->,avviso] (a1.east) -- ++(0.7,0) node[right,align=left] {仍可能直接\\连接网络};
  \node[conf,fit=(tb)(t1)] (s1) {};
  \node[above=2pt of s1] {\small （a）仅浏览器流量};
  % Tails
  \begin{scope}[xshift=4.4cm]
  \node[box] (tb2) at (0,1.6) {应用程序};
  \node[box,fill=torviola!15] (t2) at (0,-0.8) {tor};
  \node[box,fill=avviso!10] (fw) at (0,0.4) {防火墙：阻止绕过 Tor};
  \draw[->] (tb2) -- (fw);
  \draw[->] (fw) -- (t2);
  \node[conf,fit=(tb2)(fw)] (s2) {};
  \node[above=2pt of s2] {\small （b）Tails：USB Live、RAM};
  \end{scope}
  % Whonix
  \begin{scope}[xshift=8.8cm]
  \node[box] (ws) at (0,1.6) {Workstation\\（应用程序）};
  \node[box,fill=torviola!15] (gw) at (0,-0.8) {Gateway\\（tor）};
  \draw[->] (ws) -- node[right,align=left] {隔离的内部\\虚拟网络} (gw);
  \node[conf,fit=(ws)] {};
  \node[conf,fit=(gw)] {};
  \node[above=12pt of ws] {\small （c）Whonix：两台虚拟机};
  \end{scope}
  \foreach \n in {t1,t2,gw} \draw[->,thick,torverde!60!black] (\n.south) -- ++(0,-0.5) node[below] {Tor 网络};
\end{tikzpicture}
